% =============================================================
% url-href.tex — 中文 / 特殊字符 URL 写法模板
%
% 核心结论：
%   中文 / 含特殊字符的 URL 必须用 \href{raw-url}{display-text}；
%   不要用 \url{...}，也不要预先 percent-encode。
%
%   hyperref 会自动把 raw-url percent-encode 写入 PDF 注解（不会双编码）；
%   显示文本走 verbatim 解析，其中的 LaTeX 特殊字符需要转义。
%
% 使用：
%   - 单章模式下 fix.py 已自动把 \url{...} 转为 \href{raw}{display}。
%   - 本文件用于：(a) 直接撰写章节时手写 URL；(b) 维护 fix.py 时校对。
% =============================================================

% --- 正确：raw URL + raw display，hyperref 自动 percent-encode ---
% \href{https://zh.wikipedia.org/wiki/示例条目}{https://zh.wikipedia.org/wiki/示例条目}

% --- 正确：含路径中文 + 下划线 ---
% \href{https://arxiv.org/abs/1234.56789}{arXiv:1234.56789\_v2}

% --- 正确：含查询参数 + & ---
% \href{https://www.example.com/search?q=测试&type=pdf}{https://www.example.com/search?q=测试\&type=pdf}

% --- 反例（保留作对照，正常使用时删除） ---
% ❌ \url{}           ：会触发 hyperref 的 percent-decode bug，中文易出错
% ❌ \nolinkurl       ：会杀掉超链接，PDF 里只剩纯文本
% ❌ 预 percent-encoded 再进 \href：hyperref 会二次编码

% --- 显示文本必须转义的 LaTeX 特殊字符 ---
% 字符    含义       转义
% _       下划线     \_
% &       tab        \&
% $       math       \$
% #       参数       \#
% %       注释       \%
% \       命令符     \textbackslash{}

% --- 配套 hyperref 设置（main.tex / _tmp.tex 已预置，无需重复） ---
% \usepackage[unicode]{hyperref}
% \hypersetup{colorlinks=true, linkcolor=blue, urlcolor=blue, citecolor=blue}
